\chapter{Theo dõi kết quả}
\label{chuong:theo-doi}

Sau giờ mở, thầy cô xem được ai đã nộp, điểm của từng em, câu nào cả lớp sai nhiều và em nào thấy lời
giải thích của Kriky còn khó hiểu.

\viec{Bảng theo dõi}

\anh{T/28-bang-theo-doi}{
  \diem{1}{0.835}{0.284}
  \diem{2}{0.86}{0.36}
}

Bấm \nut{Bảng theo dõi} ở thanh trái. Bảng gom những việc cần thầy cô xem:

\begin{bangtra}{L{3.8cm}L{3cm}Y}{\dau{Loại} & \dau{Bấm} & \dau{Nghĩa là}}
  \nhan{Giải thích chưa rõ} & \sobuoc{1} \nut{Xem báo cáo} & Có em bấm báo Kriky giải thích khó hiểu. \\
  \nhan{Tỉ lệ sai cao} & \sobuoc{2} \nut{Xem câu} & Một câu có từ 70\% bài sai ở lần làm đầu, khi đề có trên 20 bài
  nộp (cộng các lớp). Chỉ hiện sau hạn chữa bài ở mọi lớp. \\
\end{bangtra}

Phần \nhan{Điểm gộp theo lớp} có số khi một đề đã qua hạn chữa bài ở lớp ấy. Điểm ở đây quy về thang
10.

\viec{Xem câu cả lớp sai nhiều}

\anh{T/29-cau-sai-nhieu}{
  \diem{1}{0.07}{0.885}
  \diem{2}{0.26}{0.885}
}

Màn này cho thấy mỗi phương án được bao nhiêu phần trăm học sinh chọn ở lần làm đầu, và lỗi gắn với
phương án sai. Ở ảnh, 77\% chọn \nhan{B. Cực tiểu}: phần lớn lớp nhầm chiều đổi dấu.

\begin{buoc}
  \item Xem xong bấm \nut{Đánh dấu đã xem} để câu rời khỏi bảng theo dõi.
  \item Câu có vấn đề thì bấm \nut{Chỉnh sửa trong Ngân hàng câu hỏi}.
\end{buoc}

\viec{Đọc báo cáo của học sinh}

Ở \nhan{Bảng theo dõi}, bấm \nut{Xem báo cáo} cạnh dòng \nhan{Giải thích chưa rõ}.

\anh{T/30-bao-cao}{
  \diem{1}{0.35}{0.55}
  \diem{3}{0.25}{0.92}
}

\anh[0.62\linewidth]{T/30b-bao-cao-hoc-sinh}{
  \diem{2}{0.4}{0.82}
}

\begin{buoc}
  \item Bên trái là các câu bị báo cáo. Câu cần sửa thì bấm \nut{Chỉnh sửa trong Ngân hàng câu hỏi ›}.
  \item Bên phải là từng báo cáo: em nào, sai câu nào. Bấm \nut{Mở hội thoại hỏi trợ lý ›} để đọc em
        đã hỏi Kriky những gì.
  \item Đọc xong bấm \nut{Đánh dấu đã xem}.
\end{buoc}

\anh[0.8\linewidth]{T/31-hoi-thoai-cua-em}{}

Hộp \nhan{Hội thoại hỏi trợ lý — …} chỉ để đọc. Bấm \nut{Đóng} khi xong.

\viec{Xem kết quả một đề}

\anh{T/32-cac-bai-kiem-tra}{
  \diem{1}{0.845}{0.84}
}

Ở \nhan{Các bài kiểm tra}, đề đã tới giờ mở có \nut{Xem kết quả} \sobuoc{1}. Đề phát hành cho nhiều lớp thì nút này
mở lớp đầu tiên; lớp khác mở từ màn của lớp ấy (Chương~\ref{chuong:lop}).

\anh{T/33-ket-qua-de}{
  \diem{1}{0.81}{0.11}
  \diem{2}{0.21}{0.7}
}

\begin{buoc}
  \item Bấm \nut{Xuất bảng điểm (.xlsx)} để tải bảng điểm về.
  \item Bấm tên một em để xem bài làm của em ấy.
\end{buoc}

\viec{Xem bài làm một em}

\anh{T/34-bai-lam-mot-em}{}

Màn này ghi điểm, số câu đúng ngay lần đầu, số lượt chữa đã làm và số lần em báo giải thích chưa rõ.
Bảng dưới cho thấy từng câu em chọn gì, đáp án là gì, được bao nhiêu điểm.

\viec{Hỏi Kriky về kết quả}

Thầy cô cũng hỏi Kriky được bằng lời trong một đoạn chat. Kriky trả lời số bài đã nộp, điểm trung
bình và số câu còn chữa, ví dụ khi hỏi:

\begin{loinhan}
Lớp 12A làm đề Kiểm tra 15 phút thế nào?
\end{loinhan}

Muốn biết câu nào cả lớp sai nhiều, xem mục \nhan{Tỉ lệ sai cao} của bảng theo dõi.
